\chapter{Inserting a Gastric Balloon}

\section*{Definition and Overview}
A gastric balloon is a soft, inflatable silicone balloon placed inside the stomach to treat obesity. It is inserted through the mouth and oesophagus using an endoscope, usually under sedation, and then filled with saline so that it takes up space in the stomach, making the person feel full sooner and eat less. The balloon is temporary, remaining in place for six to twelve months, and is used as a bridge to weight loss that is maintained with diet, exercise, and behavioural change. It is an option for people with obesity who have not succeeded with lifestyle measures and who are not suitable for or do not wish to have surgery. Serious complications are uncommon but include balloon deflation, blockage, and stomach ulceration.

\section*{Epidemiology}
Obesity affects around two-thirds of Australian adults, and the demand for weight-loss interventions has grown rapidly. Gastric balloons are used by a small but increasing number of people, typically those with a body mass index of 30 to 40 who have struggled with weight loss, and sometimes as a preliminary step before bariatric surgery in people who need to lose weight to reduce operative risk. Most people regain weight after balloon removal unless lasting lifestyle change is achieved, and this has shaped how the procedure is offered, with structured dietetic and psychological support as essential components.

\section*{Aetiology and Pathophysiology}
Obesity results from a sustained imbalance between energy intake and expenditure, influenced by genetics, hormones, food environment, activity, and behaviour. A gastric balloon addresses the intake side of this equation mechanically: a balloon of 400 to 700 millilitres occupies space in the stomach, distending it and activating stretch receptors that signal fullness, which delays gastric emptying and reduces the amount of food consumed at a meal. The effect is strongest in the first weeks and wanes as the stomach adapts. The balloon does not change the underlying drivers of obesity, which is why it must be combined with diet, exercise, and behaviour change to achieve lasting benefit.

\section*{Clinical Features}
\begin{itemize}
    \item The procedure is done under sedation, taking around 20 to 30 minutes
    \item Nausea, vomiting, and cramping are common in the first few days as the stomach adapts
    \item Reduced appetite and early fullness during meals
    \item Typical weight loss of around 10 to 15 per cent of body weight over six to twelve months
    \item The balloon is removed by endoscopy under sedation
    \item Mild reflux and discomfort can occur while the balloon is in place
    \item Vomiting or severe pain may indicate a complication
\end{itemize}

\section*{Differential Diagnosis}
The decision to use a gastric balloon is made after assessing the whole picture.
\begin{itemize}
    \item \textbf{Lifestyle and behavioural interventions:} appropriate as first-line treatment for most
    \item \textbf{Medicines for weight loss:} alternatives or additions to the balloon
    \item \textbf{Bariatric surgery:} more effective and durable, for people with severe obesity
    \item \textbf{Underlying causes of weight gain:} hypothyroidism, polycystic ovary syndrome, and medicines
    \item \textbf{Eating disorders:} which must be assessed before any weight-loss procedure
    \item \textbf{Medical contraindications:} stomach ulcers, previous gastric surgery, and pregnancy
\end{itemize}

\section*{When to Seek Medical Care}
People considering a gastric balloon should have a full medical and dietetic assessment and should only proceed with a structured program of support. Medical review is needed for persistent vomiting, severe abdominal pain, signs of blockage, or weight loss that stops or reverses while the balloon is in place. Any person with the balloon who develops symptoms of blockage should be assessed urgently, as a deflated balloon can obstruct the bowel.

\textbf{Call 000 (or local emergency services) for:}
\begin{itemize}
    \item Severe abdominal pain, vomiting that will not stop, or inability to keep fluids down
    \item Signs of bowel obstruction: cramping pain, vomiting, and a swollen abdomen
    \item Signs of stomach perforation: sudden severe pain, fever, and collapse
    \item Difficulty breathing or chest pain
\end{itemize}

\section*{Diagnosis and Investigations}
\begin{itemize}
    \item \textbf{Assessment of suitability:} body mass index, health conditions, medicines, and expectations
    \item \textbf{Endoscopy before placement:} to check the stomach and oesophagus for problems such as ulcers
    \item \textbf{Blood tests:} including blood glucose, lipids, and liver function
    \item \textbf{Mental health assessment:} screening for eating disorders and emotional eating
    \item \textbf{Dietetic review:} preparation for the dietary change needed
    \item \textbf{Imaging if complications are suspected:} X-ray or CT for balloon position and bowel obstruction
    \item \textbf{Regular review:} monitoring of weight, symptoms, and progress throughout
\end{itemize}

\section*{Management}
\begin{itemize}
    \item \textbf{Insertion:} the folded balloon is placed by endoscopy and filled with saline and dye
    \item \textbf{Initial symptoms:} anti-nausea medicines and a soft, small-meal diet in the first days
    \item \textbf{Structured support:} dietetic counselling, exercise guidance, and behavioural therapy
    \item \textbf{Monitoring:} regular appointments to track weight, symptoms, and nutrition
    \item \textbf{Removal:} the balloon is deflated and removed by endoscopy after six to twelve months
    \item \textbf{Transition planning:} a plan to maintain weight after removal
    \item \textbf{Treating complications:} early removal if vomiting, pain, or ulceration develops
    \item \textbf{Follow-on treatment:} medicines, surgery, or continued lifestyle programs as needed
\end{itemize}

\section*{Complications}
\begin{itemize}
    \item Nausea, vomiting, and cramping, especially early on
    \item Reflux and heartburn
    \item Ulcers and inflammation of the stomach lining
    \item Balloon deflation with migration and bowel obstruction, an emergency
    \item Obstruction of the stomach outlet by the balloon
    \item Rarely, perforation of the stomach or oesophagus
    \item Weight regain after removal without sustained lifestyle change
    \item Nutritional deficiencies if intake is too restricted
\end{itemize}

\section*{Prognosis and Outcomes}
Gastric balloons produce meaningful weight loss, averaging around 10 to 15 per cent of body weight, with improvements in blood pressure, blood sugar, and other obesity-related conditions. The effect is temporary: without lasting changes in diet and activity, much of the weight is regained within a year or two of removal. Success is therefore defined by the program around the balloon as much as by the device itself. In suitable, well-supported people, the balloon is a useful, reversible option that can also prepare people for surgery.

\section*{Prevention and Self-Care}
\begin{itemize}
    \item See the balloon as a tool within a complete lifestyle program, not a cure
    \item Follow the dietetic plan, with small, frequent, protein-rich meals
    \item Drink fluids between rather than with meals
    \item Take anti-nausea medicines as advised in the first days
    \item Attend all scheduled reviews and report symptoms promptly
    \item Build physical activity gradually
    \item Plan for life after removal before the balloon comes out
    \item Seek help early for persistent vomiting, pain, or any change in symptoms
\end{itemize}

\section*{Sources and Further Reading}
The following reputable sources provide further information for healthcare students and patients seeking reliable, current advice about gastric balloons.
\begin{itemize}
    \item Obesity Australia. \textit{Weight management information.} https://www.obesityaustralia.org/
    \item Healthdirect Australia. \textit{Weight loss procedures and devices.} https://www.healthdirect.gov.au/
    \item Gastroenterological Society of Australia. \textit{Endoscopic weight loss procedures.} https://www.gesa.org.au/
    \item Royal Australasian College of Surgeons. \textit{Bariatric treatment information.} https://www.surgeons.org/
    \item NHS. \textit{Gastric balloon.} https://www.nhs.uk/
    \item Mayo Clinic. \textit{Gastric balloon.} https://www.mayoclinic.org/
\end{itemize}
